\section*{Hoofdstuk \ref{ch:g1:my-body} --- Mijn lichaam}

\begin{solution}{exo:g1:my-body:1}
\omterm{def:g1:my-body:parts}{Hoofd}, \omterm{def:g1:my-body:parts}{romp}, \omterm{def:g1:my-body:parts}{ledematen} (twee \omterm{def:g1:my-body:parts}{armen} en twee \omterm{def:g1:my-body:parts}{benen}). De \omterm{def:g1:my-body:parts}{hals} verbindt het
\omterm{def:g1:my-body:parts}{hoofd} met de \omterm{def:g1:my-body:parts}{romp}.
\end{solution}

\begin{solution}{exo:g1:my-body:2}
Een \omterm{def:g1:my-body:parts}{arm} eindigt in een hand; een \omterm{def:g1:my-body:parts}{been} eindigt in een voet.
\end{solution}

\begin{solution}{exo:g1:my-body:3}
(a) De \omterm{def:g1:my-body:joint}{elleboog}. (b) De \omterm{def:g1:my-body:joint}{knie}. (c) De \omterm{def:g1:my-body:joint}{pols}.
\end{solution}

\begin{solution}{exo:g1:my-body:4}
De \omterm{def:g1:my-body:joint}{schouder} heeft de \omterm{def:g1:my-body:parts}{arm} dwars gedraaid, de \omterm{def:g1:my-body:joint}{elleboog} heeft hem gebogen,
en de \omterm{def:g1:my-body:joint}{pols} heeft de hand naar het oor gekanteld.
\end{solution}

\begin{solution}{exo:g1:my-body:5}
Hetzelfde: het algemene plan --- een \omterm{def:g1:my-body:parts}{hoofd}, een lichaam (\omterm{def:g1:my-body:parts}{romp}), vier
\omterm{def:g1:my-body:parts}{ledematen}. Anders: jij staat rechtop op twee \omterm{def:g1:my-body:parts}{benen} met twee vrije
\omterm{def:g1:my-body:parts}{armen}, terwijl de kat op alle vier loopt; en de kat heeft een \omterm{def:g1:animals-around-us:parts}{staart},
die jij niet hebt.
\end{solution}

\begin{solution}{exo:g1:my-body:6}
Als je alleen kijkt, schat je de bovenkant van het \omterm{def:g1:my-body:parts}{hoofd} schuin in en
komt het streepje te hoog of te laag. Het platte boek, waterpas en
tegen de muur, brengt de allerhoogste punt van het \omterm{def:g1:my-body:parts}{hoofd} recht naar de
muur: elke keer dezelfde eerlijke meting.
\end{solution}

\begin{solution}{exo:g1:my-body:7}
Ze laten zien dat je lichaam tussen de twee datums langer geworden is.
Nee --- groeien gaat te langzaam om te voelen; alleen de streepjes
maken het zichtbaar.
\end{solution}

\begin{solution}{exo:g1:my-body:8}
Open antwoord: goed lopen, gaan zitten, dingen oprapen, klappen,
schrijven --- bijna elke beweging blijkt knieën, heupen, ellebogen,
\omterm{def:g1:my-body:joint}{schouders} of polsen te gebruiken. Zonder \omterm{def:g1:my-body:joint}{gewrichten} zou het lichaam
bewegen als één stijve stok.
\end{solution}

\begin{solution}{exo:g1:my-body:9}
Dat elk lichaam normaal is op zijn eigen maat en met zijn eigen
tempo: sommige kinderen groeien vroeger, andere later, en allemaal
bereiken ze hun eigen volwassen maat. Lengtestreepjes zijn er om
jezelf te zien groeien, niet om met iemand te wedijveren.
\end{solution}

\begin{solution}{exo:g1:my-body:10}
Snijd bij de heupen (zodat de \omterm{def:g1:my-body:parts}{benen} naar voren kunnen buigen om te
zitten), bij de knieën (zodat de onderbenen omlaag kunnen hangen), en
hij zit nog beter met een snede bij de \omterm{def:g1:my-body:parts}{hals} en de \omterm{def:g1:my-body:joint}{schouders} om te
leunen en te balanceren. De sneden bootsen de \omterm{def:g1:my-body:joint}{gewrichten} van het
lichaam na: heup, \omterm{def:g1:my-body:joint}{knie}, \omterm{def:g1:my-body:parts}{hals}, \omterm{def:g1:my-body:joint}{schouder}.
\end{solution}
